% !TEX program = xelatex
% =============================================================================
% cn-legal-opinion · 律所法律意见书模板（A股IPO 场景样张）
%
% 框架依据：《公开发行证券的公司信息披露的编报规则第 12 号——公开发行证券的
% 法律意见书和律师工作报告》所列意见书应披露事项（2025 年申报文件精简改革后，
% 律师工作报告不再单独出具，相关内容并入法律意见书正文及附件）。
% 版式综合八家头部律所公开披露文书实测（吸收对照见 README）。
%
% 完整申报用法律意见书的正文事项全集（按编报规则第 12 号框架，汉坤/金杜等
% 真实目录通行为 23 项）依次为：
%   释义 / 一、本次发行上市的批准和授权 / 二、发行人本次发行上市的主体资格 /
%   三、本次发行上市的实质条件 / 四、发行人的设立 / 五、发行人的独立性 /
%   六、发起人、股东及实际控制人 / 七、发行人的股本及其演变 / 八、发行人的业务 /
%   九、关联交易及同业竞争 / 十、发行人的主要财产 / 十一、发行人的重大债权债务 /
%   十二、发行人重大资产变化及收购兼并 / 十三、发行人章程的制定与修改 /
%   十四、发行人股东会、董事会议事规则及规范运作 / 十五、发行人董事、高级管理
%   人员和核心技术人员及其变化 / 十六、发行人的税务 / 十七、发行人的环境保护和
%   产品质量、技术等标准 / 十八、发行人募集资金的运用 / 十九、发行人的业务发展
%   目标 / 二十、诉讼、仲裁或行政处罚 / 二十一、发行人招股说明书法律风险的评价 /
%   二十二、律师认为需要说明的其他问题 / 二十三、结论性意见
% 本样张为控制篇幅收录其中六项（批准授权/主体资格/实质条件/独立性/诉讼仲裁/
% 结论性意见），其余按同一层级体系增补 \section 即可。
%
% 示例样张：虚构的「北京市示例律师事务所」为「示例科技股份有限公司」（与
% cn-ipo 模板同一示例公司）科创板 IPO 出具的法律意见书。全部主体为占位
% （示例××/林××），日期与股本数据内部自洽。
%
% 编译：bash scripts/compile.sh templates/cn-legal-opinion/main.tex --preview
% 配套：supplemental.tex 为补充法律意见书（问询回复式、信函版式）样张
% =============================================================================
\documentclass[zihao=-4]{ctexart}   % 正文宋体小四
\usepackage{opinion}

% --- 文档元信息：改这里即可换成你的项目 --------------------------------------
\renewcommand{\opFirm}{北京市示例律师事务所}
\renewcommand{\opFirmEn}{EXAMPLE \& PARTNERS}
\renewcommand{\opClient}{示例科技股份有限公司}
\renewcommand{\opMatter}{首次公开发行股票并在科创板上市}
\renewcommand{\opConn}{的}                    % 题名连接字：的（金杜式）/之（竞天式）
\renewcommand{\opDocType}{法律意见书}
\renewcommand{\opRefNo}{示例（证）字〔2026〕第0001号}   % 文号（汉坤式）；置空不印
\renewcommand{\opDate}{二〇二六年七月}
\renewcommand{\opFascicle}{3-3-1}             % 申报卷册页码；非申报场景置空 {}
\renewcommand{\opMatterNo}{}                  % 左下案号（通力式）；本篇不用

\begin{document}

% --- 封面（对称式）＋目录（汉坤式）------------------------------------------
\opMakeCover
\setcounter{page}{2}      % 封面计为 3-3-1-1（方达/汉坤卷册惯例），目录起 2
\pdfbookmark[0]{目录}{toc}
\optoc

% --- 致语与引言（金杜/君合式）------------------------------------------------
\opTo{示例科技股份有限公司}

北京市示例律师事务所（以下简称"本所"）接受示例科技股份有限公司（以下简
称"发行人"或"公司"）的委托，担任发行人首次公开发行股票并在科创板上市
（以下简称"本次发行上市"）的专项法律顾问。

本所根据《中华人民共和国证券法》（以下简称"《证券法》"）、《中华人民共
和国公司法》（以下简称"《公司法》"）、《首次公开发行股票注册管理办法》
（以下简称"《注册管理办法》"）、《律师事务所从事证券法律业务管理办法》、
《律师事务所证券法律业务执业规则（试行）》《公开发行证券的公司信息披露的
编报规则第12号——公开发行证券的法律意见书和律师工作报告》（以下简称
"《编报规则第12号》"）等中国境内现行有效的法律、行政法规、规章和规范性
文件和中国证券监督管理委员会（以下简称"中国证监会"）的有关规定，按照律
师行业公认的业务标准、道德规范和勤勉尽责精神，就发行人本次发行上市事宜出
具本法律意见书。

% --- 声明事项（君合详式声明与金杜声明段的合并体例）---------------------------
\optocline{声明事项}
{\noindent\heiti\zihao{4}声明事项\par}\vspace{0.3\baselineskip}

（一）本所及经办律师依据《证券法》《律师事务所从事证券法律业务管理办法》
《律师事务所证券法律业务执业规则（试行）》等规定及本法律意见书出具之日以
前已经发生或者存在的事实，严格履行了法定职责，遵循了勤勉尽责和诚实信用原
则，进行了充分的核查验证，保证本法律意见书所认定的事实真实、准确、完整，
所发表的结论性意见合法、准确，不存在虚假记载、误导性陈述或者重大遗漏，并
承担相应法律责任。

（二）本所仅就与本次发行上市有关的法律问题发表意见，且仅根据中国境内（为
本法律意见书之目的，不包括中国香港特别行政区、中国澳门特别行政区及中国台
湾地区）现行有效的法律、法规、规章和规范性文件发表法律意见，并不依据任何
中国境外法律发表法律意见。

（三）本所不对有关会计、审计、资产评估、投资决策等非法律专业事项发表意见。
本法律意见书中对有关审计报告、验资报告、资产评估报告中某些数据和结论的引
述，并不意味着本所对该等数据、结论的真实性和准确性作出任何明示或者默示的
保证，本所并不具备核查和评价该等数据的适当资格。

（四）为出具本法律意见书，本所律师审查了发行人提供的有关文件及其复印件，
并基于发行人向本所作出的如下保证：发行人已提供了出具本法律意见书所必需的、
真实、准确、完整、有效的原始书面材料、副本材料、复印件或者口头证言，不存
在任何遗漏或者隐瞒；其所提供的副本材料或者复印件与正本材料或者原件一致。

（五）对于出具本法律意见书至关重要而又无法得到独立证据支持的事实，本所依
赖有关政府部门、发行人或者其他有关单位出具的证明文件发表法律意见。

（六）本法律意见书仅供发行人为本次发行上市之目的使用，未经本所书面同意，
不得用作任何其他目的。本所同意发行人在其为本次发行上市所制作的《招股说明
书（申报稿）》中自行引用或者按照中国证监会及上海证券交易所的审核要求引用
本法律意见书的相关内容，但发行人作上述引用时，不得因引用而导致法律上的歧
义或者曲解。

% --- 释义（金杜/方达式三列全线框表）------------------------------------------
\clearpage
\optocline{释\hspace{2\ccwd}义}
{\noindent\heiti\zihao{4}释\hspace{2\ccwd}义\par}\vspace{0.3\baselineskip}

在本法律意见书内，除非文义另有所指，下列词语具有下述含义：

\begin{definitions}
\term{本所}{北京市示例律师事务所}
\term{发行人、公司、示例科技}{示例科技股份有限公司}
\term{示例有限}{示例科技有限公司，发行人前身}
\term{本次发行上市}{发行人首次公开发行人民币普通股（A股）股票并在上海证券交易所科创板上市}
\term{保荐机构、示例证券}{示例证券股份有限公司}
\term{示例会计师}{示例会计师事务所（特殊普通合伙）}
\term{《审计报告》}{示例会计师出具的示会审字〔2026〕第0088号《审计报告》}
\term{报告期}{2023年1月1日至2025年12月31日}
\term{《公司章程》}{现行有效的《示例科技股份有限公司章程》}
\term{《公司章程（草案）》}{发行人本次发行上市后适用的《示例科技股份有限公司章程（草案）》}
\term{股东会}{发行人股东会，含整体变更为股份有限公司前的示例有限股东会}
\term{中国证监会}{中国证券监督管理委员会}
\term{上交所}{上海证券交易所}
\term{元、万元}{人民币元、万元}
\end{definitions}

% =============================================================================
% 正文（样张收录六项；全集二十三项见文首注释，按同一体系增补 \section 即可）
% =============================================================================
\clearpage

\section{本次发行上市的批准和授权}

\subsection{本次发行上市的批准}

2026年6月20日，发行人召开2026年第一次临时股东会，逐项审议通过了与本次发
行上市相关的下列议案：

\subsubsection{《关于公司申请首次公开发行股票并在科创板上市的议案》}

\subsubsection{《关于提请股东会授权董事会及其授权人士全权办理本次发行上市相关事宜的议案》}

\subsubsection{《关于公司本次发行上市前滚存未分配利润分配方案的议案》}

经本所律师核查，上述股东会的召集、召开程序、出席会议人员资格、表决程序及
表决结果均符合法律、法规、规范性文件及《公司章程》的规定，所形成的决议合
法有效。

\subsection{本次发行上市的授权}

根据上述股东会决议，股东会授权董事会及其授权人士全权办理与本次发行上市有
关的全部事宜，包括但不限于确定发行方案、选聘中介机构、制作及签署申报文件
等。该等授权的范围、程序合法有效，授权期限为自股东会审议通过之日起24个月。

\hl{综上，本所认为，发行人本次发行上市已经取得现阶段必要的批准和授权，尚
需经上交所审核同意并经中国证监会作出同意注册决定后方可实施。}

\section{发行人本次发行上市的主体资格}

\subsection{发行人系依法设立的股份有限公司}

发行人前身示例有限系于2018年5月10日设立的有限责任公司。2024年6月30日，示
例有限依法整体变更为股份有限公司，并于示例市市场监督管理局办理了变更登记，
取得统一社会信用代码为91××××××××××××××的《营业执照》。

\subsection{发行人系合法存续的股份有限公司}

经本所律师核查，截至本法律意见书出具之日，发行人系依法设立且合法存续的股
份有限公司，注册资本为36,000万元，股本总额36,000万股，不存在法律、法规、
规范性文件及《公司章程》规定的应当终止的情形。

\hl{综上，本所认为，发行人系依法设立且合法存续的股份有限公司，具备本次发
行上市的主体资格。}

\section{本次发行上市的实质条件}

发行人本次拟公开发行不超过12,000万股人民币普通股（A股），占发行后总股本
的比例不低于25\%。经本所律师逐项核查，发行人本次发行上市符合《证券法》
《注册管理办法》及《上海证券交易所科创板股票发行上市审核规则》规定的下列
实质条件：

\subsection{发行人本次发行上市符合《证券法》规定的条件}

\subsubsection{发行人具备健全且运行良好的组织机构，符合《证券法》第十二条第一款第（一）项的规定}

发行人已依法建立股东会、董事会等组织机构，不设监事会，由董事会审计委员会
行使《公司法》规定的监事会职权，并制定了相应的议事规则；各组织机构依法履
行职责、运行良好。

\subsubsection{发行人具有持续经营能力，符合《证券法》第十二条第一款第（二）项的规定}

根据《审计报告》并经本所律师核查，发行人报告期内主营业务突出、经营情况正
常，不存在影响持续经营能力的重大法律障碍。

\subsubsection{发行人最近三年财务会计报告被出具无保留意见审计报告，符合《证券法》第十二条第一款第（三）项的规定}

示例会计师已对发行人报告期各年度财务报表出具了标准无保留意见的《审计报告》。

\subsubsection{发行人及其控股股东、实际控制人最近三年不存在贪污、贿赂、侵占财产、挪用财产或者破坏社会主义市场经济秩序的刑事犯罪，符合《证券法》第十二条第一款第（四）项的规定}

经本所律师核查国家企业信用信息公示系统、中国裁判文书网等公开渠道并取得有
关主管部门出具的合规证明，发行人及其控股股东、实际控制人林××最近三年不
存在上述刑事犯罪。

\subsection{发行人本次发行上市符合《注册管理办法》规定的条件}

\subsubsection{发行人符合《注册管理办法》第十条规定的主体资格条件}

发行人系依法设立且持续经营三年以上的股份有限公司，具备健全且运行良好的组
织机构，相关机构和人员能够依法履行职责。

\subsubsection{发行人符合《注册管理办法》第十一条规定的独立性及规范运行条件}

发行人资产完整，业务及人员、财务、机构独立，与控股股东、实际控制人及其控
制的其他企业间不存在对发行人构成重大不利影响的同业竞争，不存在严重影响独
立性或者显失公平的关联交易（详见本法律意见书"四、发行人的独立性"）。

\subsubsection{发行人符合《注册管理办法》第十三条规定的合法合规条件}

发行人生产经营符合法律、行政法规的规定，符合国家产业政策；最近三年内，发
行人及其控股股东、实际控制人不存在贪污、贿赂、侵占财产、挪用财产或者破坏
社会主义市场经济秩序的刑事犯罪，不存在欺诈发行、重大信息披露违法或者其他
涉及国家安全、公共安全、生态安全、生产安全、公众健康安全等领域的重大违法
行为。

\subsection{发行人符合科创板定位}

发行人主营业务为人工智能大模型技术的研发与产业化，属于《上海证券交易所科
创板企业发行上市申报及推荐暂行规定》所列新一代信息技术领域，研发投入、发
明专利数量及营业收入增长率等指标符合科创属性评价标准。

\hl{综上，本所认为，发行人本次发行上市符合《证券法》《注册管理办法》及上
交所相关规则规定的实质条件。}

\section{发行人的独立性}

\subsection{发行人的业务独立}

发行人独立开展大模型解决方案、开放平台及智能应用业务，拥有独立完整的研发、
采购、销售体系，不存在依赖控股股东、实际控制人及其控制的其他企业开展业务
的情形。

\subsection{发行人的资产独立完整}

发行人合法拥有与经营有关的主要资产的所有权或者使用权，资产独立完整，不存
在资产、资金被控股股东、实际控制人及其控制的其他企业占用的情形。

\subsection{发行人的人员独立}

发行人的董事、高级管理人员均依照法定程序产生；总经理、财务负责人、董事会
秘书等高级管理人员未在控股股东、实际控制人及其控制的其他企业中担任除董事
以外的其他职务、未领取薪酬；财务人员未在上述企业中兼职。

\subsection{发行人的机构独立}

发行人已建立独立的组织机构，与控股股东、实际控制人及其控制的其他企业间不
存在机构混同的情形。

\subsection{发行人的财务独立}

发行人已建立独立的财务核算体系和财务管理制度，独立开设银行账户、独立纳税，
不存在与控股股东、实际控制人共用银行账户的情形。

\hl{综上，本所认为，发行人的业务、资产、人员、机构、财务独立于控股股东、
实际控制人及其控制的其他企业，具有完整的业务体系和直接面向市场独立经营的
能力。}

\section{诉讼、仲裁或行政处罚}

\subsection{发行人的重大诉讼、仲裁及行政处罚}

经本所律师核查中国裁判文书网、中国执行信息公开网、国家企业信用信息公示系
统、信用中国等公开渠道，并经发行人确认，截至2026年3月31日，发行人及其子
公司不存在尚未了结的或者可预见的对本次发行上市构成重大不利影响的诉讼、仲
裁案件，报告期内不存在重大行政处罚。

\subsection{发行人控股股东、实际控制人及董事、高级管理人员的重大诉讼、仲裁及行政处罚}

经本所律师核查上述公开渠道，并根据发行人控股股东、实际控制人及发行人董事、
高级管理人员填写的调查表，截至2026年3月31日，上述主体不存在尚未了结的或
者可预见的重大诉讼、仲裁案件，亦不存在重大行政处罚。

\section{结论性意见}

综上所述，本所认为，发行人系依法设立且合法存续的股份有限公司，具备本次发
行上市的主体资格；本次发行上市已经取得现阶段必要的批准和授权，符合《证券
法》《注册管理办法》及上交所相关规则规定的实质条件；发行人的业务、资产、
人员、机构、财务独立；发行人及相关主体不存在对本次发行上市构成重大不利影
响的诉讼、仲裁或者行政处罚。\hl{发行人本次发行上市尚需经上交所审核同意并
报经中国证监会履行发行注册程序。}

\opOriginals{肆}

% --- 签署页（金杜/方达式；「负责人」称谓各所有别：金杜「单位负责人」、
%     方达/竞天「负责人」、通力「事务所负责人」，按需改标签）----------------
\opSigPage{%
  \opsigblock{负责人：}{\opsigrow{林××}{}}
  \vspace{0.6\baselineskip}
  \opsigblock{经办律师：}{\opsigrow{王××}{张××}}}

\end{document}
